\chapter{Introduction}\label{ch:introduction}


\section{Graphene}
Graphene comprises of a single layer of carbon atoms in a honeycomb lattice. Since it's isolation in 2004 \parencite{novoselov_electric_2004}, graphene has been studied extensively due to it's unique structure and properties. Due to it's strength and conductivity, it is an ideal material for high-speed electronic devices, having applications in a range of fields, including electronics, biomedicine, photonics, etc \parencite{singh_graphene_2025}. Additionally, applications often require confinement into finite structures, such as nanoribbons and nanoflakes. \\

\noindent For an infinite, ideal graphene sheet, the conduction and valence bands meet at the Dirac points, giving a zero band gap. However, for a finite graphene flake, the gap can depend on size, shape, edge structure and symmetries \parencite{jeongCalculationElectronicStructures2024}. \\

- Compare Zig zag and armchair edges \\


\section{H\"uckel Molecular Orbital Theory}
- Nearest neighbour $\beta$, on-site energy $\alpha$ \\
- cite original method - huckel \\
- application in graphene flakes
- energy spectra, HOMO-LUMO gaps

